\documentclass[10pt,a4paper]{amsart}
\usepackage{amsmath,amssymb,amsthm,graphicx}
\usepackage{amsmath}
\usepackage{amssymb}
\usepackage{hyperref}
\usepackage[utf8]{inputenc}
\usepackage{csquotes}
\usepackage[foot]{amsaddr}


\usepackage{times}
\usepackage{amsmath,amsthm,amssymb}
\usepackage{array}
\usepackage{bbm}
\usepackage{xcolor}
\usepackage{verbatim}
\usepackage{url}
%\usepackage{authblk}
\pagestyle{plain}


\def\ve{\varepsilon}
\def\vf{\varphi}
\def\vn{\varnothing}

\def\AA{\mathcal{A}}
\def\AAA{\mathbb{A}}
\def\BB{\mathcal{B}}
\def\CC{\mathcal{C}}
\def\CCC{\mathbb{C}}
\def\DDD{\mathbb{D}}
\def\DD{\mathcal{D}}
\def\FF{\mathcal{F}}
\def\HH{\mathcal{H}}
\def\II{\mathcal{I}}
\def\III{\mathbb{I}}
\def\MM{\mathcal{M}}
\def\NN{\mathcal{N}}
\def\NNN{\mathbb{N}}
\def\PPP{\mathbb{P}}
\def\QQQ{\mathbb{Q}}
\def\RR{\mathcal{R}}
\def\RRR{\mathbb{R}}
\def\SS{\mathcal{S}}
\def\SSS{\mathbb{S}}
\def\UU{\mathcal{U}}
\def\VV{\mathcal{V}}
\def\WW{\mathcal{W}}
\def\ZZ{\mathcal{Z}}
\def\ZZZ{\mathbb{Z}}
\def\PP{\mathcal{P}}
\def\TTT{\mathbb{T}}
\def\EEE{\mathbb{E}}


%%%%  jb abbreviations
\def\bK{{\bf K}}
\def\bL{{\bf L}}
\def\bU{{\bf U}}
\def\aut{\mathop{aut}}
\def\betabar{{\bf\beta}}
\def\Uscr{\mathcal U}
\def\myreduct{\upharpoonright}
\def\abar{{\bf a}}
\def\bbar{\bf b}
\def\cbar{\bf c}
\def\lg{\mathop{\rm lg}}
\def\nmdeg{\mathop{\rm nmdeg}}
\def\icl{\mathop{\rm icl}}
% \sidebar{..} puts .. in comment form
\newcommand{\sidebar}[1]{\vskip10pt\noindent
 \hskip.70truein\vrule width2.0pt\hskip.5em
 \vbox{\hsize= 4truein\noindent\footnotesize\relax #1 }\vskip10pt\noindent}



 %%%%

\def\QQQa{\mathbb{Q}\langle \Bar{\alpha},\alpha_0 \rangle}

\def\cc{\mathfrak{c}}

\def\a{\alpha}
\def\b{\beta}
\def\g{\gamma}
\def\d{\delta}
\def\e{\varepsilon}
\def\z{\zeta}
\def\l{\lambda}
\def\k{\kappa}
\def\o{\theta}
\def\r{\rho}
\def\s{\sigma}
\def\t{\tau}
\def\w{\omega}
\def\D{\Delta}

\def\nmdeg{\operatorname{nmdeg}}
\newcommand{\Fmb}[1]{F_{Mb}\left( #1  \right)}
\newcommand{\indf}[1]{F_{ind}\left( #1  \right)}
\newcommand{\fcb}[1]{F_{cb}\left( #1  \right)}
\newcommand{\ficb}[1]{F_{Icb}\left( #1  \right)}





\DeclareMathOperator{\alg}{alg}
\DeclareMathOperator{\RU}{RU}
\DeclareMathOperator{\RM}{RM}
\DeclareMathOperator{\RD}{RD}
\DeclareMathOperator{\tp}{tp}
\DeclareMathOperator{\acl}{acl}
\DeclareMathOperator{\dcl}{dcl}
\DeclareMathOperator{\Cb}{Cb}
\DeclareMathOperator{\dtr}{\delta-trdeg}


\def\id{\operatorname{id}}
\def\Frac{\operatorname{Frac}}
\def\Const{\operatorname{Const}}
\def\spec{\operatorname{Spec}}
\def\span{\operatorname{span}}
\def\exc{\operatorname{Exc}}
\def\Div{\operatorname{Div}}
\def\cl{\operatorname{cl}}
\def\mer{\operatorname{mer}}
\def\trdeg{\operatorname{trdeg}}
\def\ord{\operatorname{ord}}




\newtheorem{theorem}{Theorem}[section]
\newtheorem{thm}[theorem]{Theorem}
\newtheorem{question}[theorem]{Question}
\newtheorem{lemma}[theorem]{Lemma}
\newtheorem{cor}[theorem]{Corollary}
\newtheorem{prop}[theorem]{Proposition}
\newtheorem{fact}[theorem]{Fact}
\newtheorem{claim}[theorem]{Claim}
\newtheorem{subclaim}[theorem]{Subclaim}
\newtheorem{conj}[theorem]{Conjecture}
\newtheorem{cond}[theorem]{Condition}
\newtheorem{observation}[theorem]{Observation}
\newtheorem{notation}[theorem]{Notation}

\theoremstyle{definition}
\newtheorem{definition}[theorem]{Definition}
\newtheorem{example}[theorem]{Example}
\newtheorem{xca}[theorem]{Exercise}
\newtheorem{ques}[theorem]{Question}
\newtheorem{note}[theorem]{Note}

\theoremstyle{remark}
\newtheorem{remark}[theorem]{Remark}
\newcommand{\m}{\mathbb }
\newcommand{\mc}{\mathcal }
\newcommand{\mf}{\mathfrak }
\newcommand{\mb}{F_{Mb} }
\newcommand{\G}{\mathrm{G} }
\newcommand{\U}{\mathrm{U} }



\def \ta {\tau_{\mathcal{D}/\Delta}}
\def \D {\Delta}
\def \DD {\mathcal D}

\newcommand{\gen}[1]{\left\langle#1\right\rangle}
%\renewcommand{\restriction}{\mathord{\upharpoonright}}

\newcommand{\pd}[2]{\frac{\partial #1}{\partial #2}}
\newcommand{\td}{\text{tr.deg.}}
\newcommand{\pp}{\partial }
\newcommand{\pdtwo}[2]{\frac{\partial^2 #1}{\partial #2^2}}
\newcommand{\od}[2]{\frac{d #1}{d #2}}
\def\Ind{\setbox0=\hbox{$x$}\kern\wd0\hbox to 0pt{\hss$\mid$\hss} \lower.9\ht0\hbox to 0pt{\hss$\smile$\hss}\kern\wd0}
\def\Notind{\setbox0=\hbox{$x$}\kern\wd0\hbox to 0pt{\mathchardef \nn=12854\hss$\nn$\kern1.4\wd0\hss}\hbox to 0pt{\hss$\mid$\hss}\lower.9\ht0 \hbox to 0pt{\hss$\smile$\hss}\kern\wd0}
\def\ind{\mathop{\mathpalette\Ind{}}}
\def\nind{\mathop{\mathpalette\Notind{}}}
\numberwithin{equation}{section}
\newcommand{\is}{^{p^ {-\infty}}}
\newcommand{\QQ}{\mathbb Q}
\newcommand{\fh}{\mathfrak h}
\newcommand{\SL}{\operatorname{SL}}


\title{Simple Homogeneous Structures and Indiscernible Sequence Invariants}
\author[1]{John T. Baldwin}
\author[2]{James Freitag}
\author[3]{Scott Mutchnik}
%\affil[1]{Department of Mathematics, Statistics, and Computer Science, University of Illinois at Chicago}
%%%\affil[2]{Department of Biology, University Y}
%%\author{ John T. Baldwin}
%\author{James Freitag} 
\address{Department of Mathematics, Statistics, and Computer Science, University of Illinois at Chicago }
\email{jbaldwin@uic.edu}
\email{jfreitag@uic.edu}

%\author{Scott Mutchnik}
%\address{Department of Mathematics, Statistics, and Computer Science, University of Illinois at Chicago}
\email{mutchnik@uic.edu}

\thanks{JB was partially supported by Simons travel grant G3535. JF was partially supported by NSF CAREER award number 1945251. SM, who is additionally affiliated with the Institut de mathématiques de Jussieu – Paris Rive Gauche, was supported by the NSF under Grant No. DMS-2303034, and by the European Union’s Horizon 2020 research and innovation programme under the Marie Sklodowska-Curie grant agreement No 101034255.}

\date{\today}
\title{Non-minimality degree can be made arbitrarily large in an omega-categorical supersimple theory}
\author{John Baldwin, James Freitag, Scott Mutchnik}
\date{arxiv:2405.08211v3; CC BY 4.0}
\begin{document}
\maketitle
\noindent\textit{Editorial note.} All delivered mathematical text is the author's own, copied from the pinned publisher source sources/logic-fl03-simple-homogeneous-structures-and/pinned-original-source.tex; no statement or proof step was rewritten, reordered or removed.\par\medskip
\noindent\textit{Editorial note.} Preamble provenance: the retained preamble is the source's own layout adapter, if any, plus the author's macro block at lines 3-205 of sources/logic-fl03-simple-homogeneous-structures-and/pinned-original-source.tex (everything between the publisher class metadata and \textbackslash{}begin\{document\}, comments excluded). Nothing was added, no macro was invented, and the publisher document class is deliberately not reproduced; the excerpt serves only as preamble, never as a span coordinate.\par\medskip
\noindent\textit{Editorial note.} Metadata provenance: source\_id and source\_version were taken from sources/logic-fl03-simple-homogeneous-structures-and/license-observation.json; the title from SOURCE.json; the author list from the source metadata record; the source URL from the first non-empty of source\_url, source\_urls[0] and the evidence's fixed\_version\_url/license\_url. The licence evidence file used by this card is sources/logic-fl03-simple-homogeneous-structures-and/license-observation.json, recorded in license.evidence\_path so that the binding target is explicit.\par\medskip
\noindent\textit{Editorial note.} External references are navigation pointers only; every deliberately kept step of the author's argument stays inside the counted spans. The displayed numbers of those pointers are the author's own numbering under the author's \textbackslash{}newtheorem declarations (section.index as declared), computed mechanically by the lane's scanner.\par\medskip
\noindent\textit{Editorial note.} The counted proof contains no reference to any labelled item and no bibliography citation, so the card needs neither a navigation stub nor a bibliography excerpt.\par\medskip
\noindent\textit{Editorial note.} The statement is staged as the author's 'fact' environment; its span includes the author's own label line placed after the environment body, which is why the label sits at source line 1722.\par\medskip
\noindent\textit{Editorial note.} Assembly history: r001 failed with 'Option clash for package graphics' because the author's preamble loads \textbackslash{}usepackage[dvips]\{graphics,epsfig\} on top of the wrapper class's graphicx baseline; r002 splits the preamble and omits that single package line. The generator now performs this split automatically, so later cards from this source will not hit it again. No mathematics changed.\par\medskip
\noindent\textit{Allowed standard prerequisites (editorial).}\begin{itemize}
\item Evans's construction of *-generic structures from a class C\_f with a predimension, and his characterization of forking by predimension. Used inside the counted proof at lines 1726-1733. Supplies the theory T\_n and the criterion that non-forking is detected by the equality d(A/C) = d(A/BC).
\item The criteria on f for the generic theory to be well-defined and simple. Verified inside the counted proof at lines 1732-1733. Shows that the chosen f(x) = log(x)/log(3) satisfies 0 < f' < 1/x and f(3x) <= f(x)+1, so that T\_n is supersimple.
\end{itemize}\medskip
\begin{fact}
    For any $n < \omega$, there is an $\omega$-categorical supersimple theory $T_n$
    admitting a type $p$ with $\mathrm{nmdeg}(p) \geq n$.
\label{simpnm}
\end{fact}

\begin{proof} Let $\mathcal{L} = \{R\}$ for $R$ a
$k$-ary relation  with $k = 3^{2n+1}$, and let $\delta(A) = 2|A| - \frac{|R(A)|}{k!}$. (Note that the factor $\frac{1}{k!}$ accounts for ths symmetry, so $\frac{|R(A)|}{k!}$ is equal to the number of tuples with coordinates in $A$, up to permutation, that realize $R(A)$).  Then let $T_n$ be the theory of the *-generic $M$ for the class $\mathcal{C}_{f}$,
where $f(x) = \frac{\mathrm{log}(x)}{\mathrm{log}(3)}$.




Now the criteria for $T$ being well-defined and simple, $0 < f'(x) <
\frac{1}{x}$ and $f(3x) \leq f(x) + 1$, will hold for all $x \geq 0$. By Evans's characterization of forking ($A \ind _{C} B$ if and only if $d(A/C) = d(A/BC)$ for $A$, $B$, $C$ finite sets), and the choice of $d_{0}$, $T_{n}$ is supersimple of $\mathrm{SU}$-rank at most $2$. It remains to show that there is a unique type in $S^{n}(\emptyset)$ realized by an $n$-tuple of distinct elements (i.e. $T_{n}$ is $n$-transitive), and that for $p \in S^{1}(\emptyset)$ the unique $1$-type, $\mathrm{SU}(p) = 2$. To show the first claim, let $a_{1}, \ldots, a_{n} \in \mathbb{M}$ be distinct; it suffices to show that $\{a_{1}, \ldots, a_{n}\}$ is algebraically closed, because isomorphic algebraically closed sets have the same type. But $A=  \mathrm{acl}(a_{1}, \ldots, a_{n})$ must have predimension at most $d_{0}(\{a_{1}, \ldots a_{n}\}) = 2n$. So if $\{a_{1}, \ldots, a_{n}\}$ is not algebraically closed, so $|A| > n$, it must then be the case that $|A| \geq k$, because otherwise $A$ would have no instances of $R$ and $d_{0}(A)= 2|A| > 2n$. But then by choice of $f$ and $k$, $f(|A|) > 2n$, contradicting $d_{0}(A) \leq 2n$.

To show $SU(p) \geq 2$, so $\mathrm{SU}(p) = 2$, because $f(x) \leq 2x -1$ for $x \leq k$, one can see there is a structure in  $\mathcal{C}_{f}$ of size $k$ with exactly one instance (up to permutation) of the relation $R$. Since $T_{n}$ is the theory of the $*$-generic of $\mathcal{C}_{f}$, there is an algebraically closed set $\{a_{1}, \ldots, a_{k}\} \subseteq \mathbb{M}$ isomorphic to this structure. Since $\{a_{1}, \ldots, a_{k}\}$ and therefore $\{a_{1}\}$, $\{a_{2} \ldots, a_{k}\}$ and $\emptyset$ are algebraically closed, $d(a_{1}/a_{2} \ldots a_{k}) = d_{0}(a_{1} \ldots a_{k}) - d_{0}(a_{2} \ldots a_{k})= 1$ and $a_{1} \notin \mathrm{acl}(a_{2}, \ldots a_{k})$, while $d(a_{1}/\emptyset)=d_{0}(a_{1}) - d_{0}(\emptyset) = 2$, so $a_{1} \nind_{\emptyset} a_{2} \ldots a_{k}$. Since $a_{1} \models p$, $\mathrm{SU}(p) \geq 2$, so $\mathrm{nmdeg}(p) \geq n$ by the $n$-transitivity shown above.


\end{proof}

\end{document}
